\documentclass[tikz,border=4mm]{standalone}

\usepackage{fontspec}
\setmainfont{LibertinusSans-Regular.otf}[
  BoldFont=LibertinusSans-Bold.otf,
  ItalicFont=LibertinusSans-Italic.otf,
  BoldItalicFont=LibertinusSans-Bold.otf,
]

\usepackage{xcolor}
\usetikzlibrary{arrows.meta, calc, positioning, fit, backgrounds}

\definecolor{Ink}{HTML}{243042}
\definecolor{Soft}{HTML}{F7F9FC}
\definecolor{Panel}{HTML}{EEF4FB}
\definecolor{AccentA}{HTML}{246BCE}
\definecolor{AccentB}{HTML}{C8463A}
\definecolor{Muted}{HTML}{5D6B7A}

\tikzset{
  every node/.style={text=Ink},
  block/.style={
    draw=Ink!85,
    line width=0.5pt,
    rounded corners=3pt,
    fill=white,
    minimum width=31mm,
    minimum height=9mm,
    inner sep=3pt,
    align=center,
    font=\scriptsize
  },
  core/.style={
    block,
    minimum height=10mm,
    fill=Panel
  },
  sidecar/.style={
    draw=Ink!65,
    line width=0.45pt,
    rounded corners=3pt,
    fill=Soft,
    minimum width=19mm,
    minimum height=7mm,
    inner sep=2.5pt,
    align=center,
    font=\tiny
  },
  flowa/.style={-Latex, line width=1.5pt, draw=AccentA},
  flowb/.style={-Latex, line width=1.5pt, draw=AccentB},
  lookup/.style={<->, line width=0.8pt, draw=Ink!65, densely dashed},
  pilla/.style={
    fill=AccentA!9,
    draw=AccentA!45,
    rounded corners=6pt,
    inner xsep=4pt,
    inner ysep=2pt,
    text=AccentA!80!black,
    font=\bfseries\tiny
  },
  pillb/.style={
    fill=AccentB!9,
    draw=AccentB!45,
    rounded corners=6pt,
    inner xsep=4pt,
    inner ysep=2pt,
    text=AccentB!80!black,
    font=\bfseries\tiny
  },
  note/.style={
    text width=31mm,
    align=left,
    font=\tiny,
    text=Muted
  }
}

\begin{document}
% Replace labels and tune spacing, but preserve the general top-down structure.
% This template is optimized for marginalia, not for wide figures in the main text.
\begin{tikzpicture}[x=1cm,y=1cm]
  \node[block] (step1) at (0,0) {Шаг 1};
  \node[block, below=7.5mm of step1] (step2) {Шаг 2};
  \node[block, below=7.5mm of step2] (step3) {Шаг 3};
  \node[core, below=7.5mm of step3] (core) {Ключевой узел};
  \node[block, below=10.5mm of core] (step5) {Шаг 5};

  \begin{scope}[on background layer]
    \node[
      draw=Ink!10,
      fill=Soft,
      rounded corners=6pt,
      fit=(step1) (step2) (step3) (core) (step5),
      inner xsep=9pt,
      inner ysep=10pt
    ] {};
  \end{scope}

  \draw[flowa] (step1.south) -- (core.north);
  \draw[flowb] (core.south) -- (step5.north);

  \node[pilla, anchor=west] at ($(step1.east)+(4mm,-1mm)$) {Поток A};
  \node[pillb, anchor=west] at ($(core.east)+(4mm,-1mm)$) {Поток B};

  \node[sidecar, right=7mm of core] (sidecar) {Служебный\\узел};
  \draw[lookup] (core.east) -- (sidecar.west);
  \node[font=\tiny, text=Muted, anchor=north west] at ($(sidecar.south west)+(0,-1.5mm)$) {lookup A $\rightarrow$ B};

  \node[note, below=7mm of step5] {Короткое пояснение под схемой. Если пояснение не помещается без тесноты, это уже не маргинальная схема.};
\end{tikzpicture}
\end{document}
